\section{Discussion}
\label{sec:discussion}

\paragraph{What the evidence supports.}
Bridge failure is the dominant source of loss on bridge questions: for every generator, answer F1 collapses and abstention rises sharply when the context lacks a gold paragraph (Table~\ref{tab:link}). Closing the retrieval gap is therefore worth a large amount of answer quality, up to the 71--79\% F1 of the gold-first oracle. Among methods that make no LLM call, SES-RAG has the highest gold-complete rate at all five budgets we tested (300 to 1{,}500 tokens), by a margin that shrinks to under one point at 1{,}500 tokens (Fig.~\ref{fig:budget}); at 300 tokens it is nominally ahead of the iterative baseline as well (47.3\% vs.\ 46.3\%, not tested). Its F1 advantage over Hybrid RAG is stable across four generators. Its F1 lead over GraphRAG-lite and SpectRAG-lite is 1.2 to 3.7 points; the uncorrected intervals exclude zero for three of four generators against GraphRAG-lite and two of four against SpectRAG-lite, but not under the Holm correction used for the primary generator.

\paragraph{What it does not support.}
Three claims that one might expect from a spectral formulation are not supported. (i) SES-RAG is not significantly better than the strongest graph baselines (GraphRAG-lite, SpectRAG-lite) after Holm correction. (ii) $\lambda_2$ is not the active ingredient: replacing it by the number of connected components (a simpler quantity) or by an unnormalised $\lambda_2$ matches or exceeds it (Table~\ref{tab:abl}), so what matters is rewarding a \emph{connected} context, not the particular spectral measure. The structure of the normalised $\lambda_2$ offers a likely explanation: it falls as a connected set grows, so the reward mostly penalises fragmentation and gives no credit for extending a chain (Section~\ref{sec:method}); a component count measures fragmentation directly and is unaffected by growth. (iii) An LLM-driven iterative retriever is clearly ahead, by 11 points of gold-complete rate and 7--10 points of F1 depending on the generator, at the price of three LLM calls per question.

\paragraph{Where the loss occurs.}
The candidate pool contains all gold paragraphs for 90.8\% of questions, yet SES-RAG's context is gold-complete for 63.1\%. Pool recall is a loose ceiling, since the pool holds 100 candidates for about five context slots, but it shows that the relevance score and the connectivity reward often fail to single out the bridge among many candidates. The iterative retriever, which finds the bridge with LLM-written queries, is 11 points ahead, and packing its candidates with our rule does not help. Connectivity alone does not discriminate: in 26.7\% of questions the two gold paragraphs have no direct graph edge, and distractors connected to the seeds compete for the same budget.

\paragraph{Dependence on the graph.}
The signal that makes connectivity informative is the title-mention edge, which mirrors hyperlinks between encyclopedia articles. Without it, the advantage of SES-RAG over the graph baselines vanishes (Table~\ref{tab:nomention}). This is consistent with the idea that a connectivity reward can only help if the graph encodes the bridge relation; we did not test this with other kinds of links. The mention edges mirror the hyperlink structure that defines HotpotQA's bridges, so the result is also a statement about how favourable the benchmark graph is. Corpora without explicit links would need link prediction, entity linking or an LLM-built graph.

\paragraph{Practical guidance.}
When query-time LLM calls are acceptable, iterative retrieval is the stronger choice. When latency or cost forbids them, a connectivity-aware selector over a link-aware graph is the best option we evaluated, and selection takes about four milliseconds per question on a CPU. Combining the two (IRCoT+SES) did not help on the primary set and lowered gold-complete retrieval on MuSiQue, which suggests that the benefit of connectivity-aware packing is small once the retriever already finds the bridge.

\section{Limitations}
\label{sec:limits}
\begin{itemize}
\item \textbf{Simplified baselines.} GraphRAG-lite, SpectRAG-lite, G-Retriever-lite and IRCoT-lite re-implement the retrieval principle of each system on one shared graph; they omit LLM-built graphs, community summaries, learned components and, for IRCoT, chain-of-thought reasoning inside the loop. Conclusions about the original systems cannot be drawn from them.
\item \textbf{Graph construction.} Title-mention edges rely on encyclopedic titles. The corpus pools 2{,}000 questions' paragraphs, which is much smaller than a full Wikipedia index.
\item \textbf{Evaluation.} The primary dataset is HotpotQA bridge questions; the external datasets use settings frozen on HotpotQA. Against a second judge (gpt-4o, 10\% random subsample per generator) the correctness judgments agree 91.6--93.4\% ($\kappa=0.62$--$0.75$), but the faithfulness judgments only 83.6--88.3\% ($\kappa=0.18$--$0.53$), so we do not base any claim on faithfulness or the unsupported-answer rate. Both judges are LLMs and no human validation was done; claude-haiku-5-5 judges its own answers as well as those of the other generators.
\item \textbf{Generators.} Four generators from two vendors were used; the Claude models accept no temperature parameter and ran with default sampling, so their results are not strictly deterministic. A third vendor could not be evaluated at scale because of API quota limits.
\item \textbf{Scope.} Comparison questions, which need parallel rather than chained evidence, are outside the study. In an earlier mixed-type run (different generator and protocol, so only indicative) SES-RAG retrieved a gold-complete context for 59\% of comparison questions against 89--95\% for the baselines, because the connectivity reward pulls the selection towards a chain.
\item \textbf{External-dataset answer metrics.} On 2WikiMultiHopQA and MuSiQue the gold-first oracle abstains on about 48\% of questions, so F1 is capped far below HotpotQA's and we rely on the gold-complete rate there.
\item \textbf{Answer format and judging.} The Claude generators produce answers of 60 or more characters for about 8\% of questions against under 1\% for the OpenAI models, so F1 and judged accuracy are not format-controlled across vendors, and Sonnet and Haiku were run with different thinking settings (necessarily, as the API requires different values). The judge's accuracy exceeds the F1-threshold accuracy by 2.5 points for both Claude generators and by 0.9 and 1.6 points for gpt-4o-mini and gpt-4o.
\item \textbf{Statistics.} One generation run per generator and system; no human validation of the judge; ablation intervals are not corrected beyond the primary family; the $\mu$ and budget sweeps were run on the primary set and are exploratory.
\item \textbf{Artifacts.} The code and outputs are in the project repository; a public repository URL and archival DOI are to be added at submission.
\end{itemize}
